\documentclass[11pt,a4paper]{article}

\usepackage[T1]{fontenc}
\usepackage[utf8]{inputenc}
\usepackage{lmodern}
\usepackage[margin=2.2cm]{geometry}
\usepackage{amsmath,amssymb,amsthm,mathtools}
\usepackage{bm}
\usepackage{booktabs}
\usepackage{array}
\usepackage{xcolor}
\usepackage{tikz}
\usetikzlibrary{arrows.meta,calc,positioning}
\usepackage{enumitem}
\usepackage[colorlinks=true,linkcolor=blue!60!black,citecolor=blue!60!black,urlcolor=blue!60!black]{hyperref}

\setlength{\emergencystretch}{3em}
\setlength{\parskip}{4pt}

\definecolor{atomblue}{RGB}{41,128,185}
\definecolor{defectred}{RGB}{211,47,47}
\definecolor{okgreen}{RGB}{46,125,50}
\definecolor{boxbg}{RGB}{245,248,252}
\definecolor{boxframe}{RGB}{180,200,225}

\newcommand{\R}{\mathbb{R}}
\newcommand{\Z}{\mathbb{Z}}
\newcommand{\Q}{\mathbb{Q}}
\newcommand{\T}{^{\mathsf T}}
\newcommand{\code}[1]{\texttt{#1}}

\newenvironment{keyidea}[1][Key Physical \& Mathematical Principle]{%
  \par\medskip\noindent
  \begin{tikzpicture}
    \node[draw=boxframe,fill=boxbg,rounded corners=4pt,inner sep=10pt,text width=\dimexpr\textwidth-22pt\relax] (box) \bgroup
    \textbf{\textcolor{blue!60!black}{#1.}}\par\smallskip
}{%
    \egroup;
  \end{tikzpicture}\par\medskip
}

\title{\textbf{Mathematical Formulation of Special Grain Boundaries and Associated Boundary Conditions in the Mesoscopic Tensorial Model (MTM)}\\[6pt]
\large Coincidence Site Lattice (CSL) Rationalization, Seam Overlap Resolution, Analytical Pre-Opening, Energy-Monotone Delaunay Surgery, and Two-Stage Relaxation--Loading Boundary Algebra}
\author{MTM Lattice-Triangulation Documentation}
\date{November 2026}

\begin{document}
\maketitle

\begin{abstract}
\noindent
This document presents the complete mathematical formulation and algorithmic architecture for constructing special Coincidence Site Lattice (CSL) grain boundaries and enforcing their associated two-stage boundary conditions (\textbf{Relaxation Step BCs} vs.\ \textbf{Loading Step BCs \& Indenter}) in the 2D Mesoscopic Tensorial Model (MTM). We derive the coprime Diophantine generators $(m,n)$ for square ($C_4$) and triangular ($C_6$) Bravais lattices, the parity halving rule for the coincidence index $\Sigma(m,n)$, the structural repeat period $L_{\text{CSL}}(m,n)$, the interplanar layer spacing $d_{mn} = a_0^2 / L_{\text{CSL}}$, and the closed-form analytical normal pre-opening gap $\delta_\perp(m,n)$ that prevents initial cross-boundary triangles from starting inside the repulsive core of the Zanzotto/Ericksen potential. We then formulate the 5-state independent $(x,y)$ boundary condition algebra (\code{node\_bc\_role} $\in \{0,1,2,3,4\}$), the orthogonal subspace projection operator $\mathbf{P}_{\text{BC}}$, the proof that the 3-Roller + 1-Free \code{GB-Frame Roller} configuration eliminates all three 2D rigid-body zero modes while guaranteeing $\sigma_{\perp\perp}^{\text{eq}} = 0$, and the reference handover operator ($\mathbf{X}_0 \leftarrow \mathbf{x}_{\text{eq}}$) that eliminates boundary snap upon transitioning from Stage 1 relaxation to Stage 2 mechanical loading or nanoindentation.
\end{abstract}

\tableofcontents

% ============================================================================
\section{2D Crystallographic Kinematics and Rotated Grain Lattices}
% ============================================================================

\subsection{Bravais Lattice Bases and Point-Group Symmetry}
Let $a_0 > 0$ denote the fundamental equilibrium lattice parameter ($a_0 = 1$ in normalized lattice units). A 2D Bravais lattice $\mathcal{L}_0 \subset \R^2$ is generated by integer linear combinations of two primitive basis vectors $\mathbf{a}_1, \mathbf{a}_2 \in \R^2$:
\begin{equation}
  \mathcal{L}_0 = \left\{ \mathbf{X}_{i,j} = i\,\mathbf{a}_1 + j\,\mathbf{a}_2 \;\middle|\; (i,j) \in \Z^2 \right\} = \mathbf{A}_0\,\Z^2, \qquad \mathbf{A}_0 = \begin{bmatrix} \mathbf{a}_1 & \mathbf{a}_2 \end{bmatrix}.
\end{equation}
Two fundamental 2D crystallographic symmetry classes are implemented:
\begin{enumerate}
  \item \textbf{Square Lattice ($C_4 / D_4$ symmetry, \code{square})}:
  \begin{equation}
    \mathbf{a}_1 = a_0 \begin{bmatrix} 1 \\ 0 \end{bmatrix}, \qquad
    \mathbf{a}_2 = a_0 \begin{bmatrix} 0 \\ 1 \end{bmatrix}, \qquad
    \mathbf{A}_0^{\text{sq}} = a_0 \begin{bmatrix} 1 & 0 \\ 0 & 1 \end{bmatrix}, \qquad
    A_{\text{cell}}^{\text{sq}} = \det\mathbf{A}_0^{\text{sq}} = a_0^2.
  \end{equation}
  By $D_4$ dihedral symmetry, all physically distinct symmetric tilt misorientations lie in $\Delta\theta \in [0^\circ, 90^\circ]$, with mirror duality across $\Delta\theta = 45^\circ$.

  \item \textbf{Triangular / Hexagonal Lattice ($C_6 / D_6$ symmetry, \code{triangular})}:
  \begin{equation}
    \mathbf{a}_1 = a_0 \begin{bmatrix} 1 \\ 0 \end{bmatrix}, \qquad
    \mathbf{a}_2 = a_0 \begin{bmatrix} \tfrac{1}{2} \\ \tfrac{\sqrt{3}}{2} \end{bmatrix}, \qquad
    \mathbf{A}_0^{\text{tri}} = a_0 \begin{bmatrix} 1 & \tfrac{1}{2} \\ 0 & \tfrac{\sqrt{3}}{2} \end{bmatrix}, \qquad
    A_{\text{cell}}^{\text{tri}} = \det\mathbf{A}_0^{\text{tri}} = \frac{\sqrt{3}}{2}a_0^2.
  \end{equation}
  By $D_6$ dihedral symmetry, all physically distinct symmetric tilt misorientations lie in $\Delta\theta \in [0^\circ, 60^\circ]$, with duality across $\Delta\theta = 30^\circ$.
\end{enumerate}

For a grain $g$ with crystallographic orientation $\theta_g \in \R$ and origin anchor $\mathbf{p}_g \in \R^2$, the rotated lattice basis is $\mathbf{A}_g = \mathbf{R}(\theta_g)\mathbf{A}_0$, where
\begin{equation}
  \mathbf{R}(\theta_g) =
  \begin{bmatrix}
    \cos\theta_g & -\sin\theta_g \\
    \sin\theta_g & \cos\theta_g
  \end{bmatrix} \in \mathrm{SO}(2), \qquad
  \mathcal{L}_g(\theta_g, \mathbf{p}_g) = \left\{ \mathbf{p}_g + \mathbf{R}(\theta_g)\mathbf{A}_0\begin{bmatrix} i \\ j \end{bmatrix} \;\middle|\; (i,j) \in \Z^2 \right\}.
\end{equation}

\subsection{Bicrystal Grain Boundary Plane Geometry}
In a macroscopic domain $\Omega = [x_{\min}, x_{\min}+L_x] \times [y_{\min}, y_{\min}+L_y]$, a planar bicrystal interface $\Gamma_{\text{GB}}$ is specified by normalized pivot coordinates $(c_x, c_y) \in [0,1]^2$ and normal inclination angle $\phi_{\text{GB}}$ (\code{gb\_angle\_deg}):
\begin{equation}
  \mathbf{p}_{\text{GB}} = \begin{bmatrix} x_{\min} + c_x L_x \\ y_{\min} + c_y L_y \end{bmatrix}, \qquad
  \hat{\mathbf{n}}_{\text{GB}}(\phi_{\text{GB}}) = \begin{bmatrix} \cos\phi_{\text{GB}} \\ \sin\phi_{\text{GB}} \end{bmatrix}, \qquad
  \hat{\mathbf{t}}_{\text{GB}}(\phi_{\text{GB}}) = \begin{bmatrix} -\sin\phi_{\text{GB}} \\ \cos\phi_{\text{GB}} \end{bmatrix}.
\end{equation}
The signed normal distance $s(\mathbf{x})$ and tangential coordinate $\tau(\mathbf{x})$ of any point $\mathbf{x} \in \Omega$ are
\begin{equation}
  s(\mathbf{x}) = (\mathbf{x} - \mathbf{p}_{\text{GB}})\cdot\hat{\mathbf{n}}_{\text{GB}}, \qquad
  \tau(\mathbf{x}) = (\mathbf{x} - \mathbf{p}_{\text{GB}})\cdot\hat{\mathbf{t}}_{\text{GB}}.
\end{equation}
For a symmetric tilt grain boundary (STGB) with half-tilt angle $\theta$ and total misorientation $\Delta\theta = 2\theta$, the two half-crystals have orientations
\begin{equation}
  \theta_1 = \phi_{\text{GB}} - \theta = \phi_{\text{GB}} - \frac{\Delta\theta}{2} \quad (s \le 0), \qquad
  \theta_2 = \phi_{\text{GB}} + \theta = \phi_{\text{GB}} + \frac{\Delta\theta}{2} \quad (s \ge 0).
\end{equation}

% ============================================================================
\section{Coincidence Site Lattice (CSL) Theory and General $(m,n)$ Rationalizer}
% ============================================================================

\subsection{Diophantine Coincidence Condition on a 2D Square Lattice}
Consider two square lattices $\mathcal{L}_1 = \mathbf{R}(-\theta)\Z^2 a_0$ and $\mathcal{L}_2 = \mathbf{R}(+\theta)\Z^2 a_0$ meeting along a vertical boundary plane ($\phi_{\text{GB}} = 0^\circ$, $\hat{\mathbf{t}}_{\text{GB}} = \hat{\mathbf{e}}_y$). A lattice site $\mathbf{p} = a_0(n\hat{\mathbf{e}}_x + m\hat{\mathbf{e}}_y) \in \Z^2 a_0$ lies strictly on the vertical boundary axis after rotation by $+\theta$ if and only if its $x$-component vanishes:
\begin{equation}
  \mathbf{R}(\theta)\begin{bmatrix} n a_0 \\ m a_0 \end{bmatrix}
  = a_0 \begin{bmatrix} n\cos\theta - m\sin\theta \\ n\sin\theta + m\cos\theta \end{bmatrix}
  = \begin{bmatrix} 0 \\ L_{\text{CSL}} \end{bmatrix}.
\end{equation}
This yields the coprime integer generator law ($m, n \in \Z_{>0}$, $\gcd(m,n) = 1$, $1 \le n \le m$):
\begin{equation}
  \tan\theta_{mn} = \frac{n}{m}, \qquad
  \theta_{mn} = \arctan\!\left(\frac{n}{m}\right), \qquad
  \Delta\theta_{mn} = 2\arctan\!\left(\frac{n}{m}\right).
\end{equation}
The relative rotation matrix between the two grains is then a rational orthogonal matrix:
\begin{equation}
  \mathbf{R}(\Delta\theta_{mn}) = \frac{1}{m^2 + n^2}
  \begin{bmatrix}
    m^2 - n^2 & -2mn \\
    2mn & m^2 - n^2
  \end{bmatrix} \in \mathrm{SO}(2,\Q).
\end{equation}

\subsection{The Coincidence Index $\Sigma(m,n)$ and Parity Halving Rule}
The reciprocal density of shared coincidence sites on the square lattice is given by
\begin{equation}
  \Sigma_{\text{sq}}(m,n) = \frac{m^2 + n^2}{\gcd(m^2 + n^2,\, 2)}.
\end{equation}
\begin{keyidea}[Parity Halving and $90^\circ$ Duality of Square-Lattice CSL Boundaries]
Because $\gcd(m,n) = 1$, $m$ and $n$ cannot both be even:
\begin{itemize}[leftmargin=1.5em]
  \item \textbf{Opposite parity ($m+n$ odd $\implies m^2+n^2$ odd)}: $\gcd(m^2+n^2, 2) = 1$, so $\Sigma(m,n) = m^2 + n^2$. For $(m,n) = (2,1)$, $\Sigma = 5$ at $\Delta\theta = 2\arctan(1/2) \approx 53.130^\circ$ ($\Sigma 5\,(210)$).
  \item \textbf{Both odd ($m, n$ odd $\implies m^2+n^2$ even)}: $(m\pm n)$ are both even, so the centered vector $\frac{1}{2}((m+n)\mathbf{a}_1 + (m-n)\mathbf{a}_2) \in \Z^2 a_0$ halves the CSL unit cell area, yielding $\Sigma(m,n) = (m^2+n^2)/2$. For $(m,n) = (3,1)$, $\Sigma = 10/2 = 5$ at $\Delta\theta = 2\arctan(1/3) \approx 36.870^\circ$ ($\Sigma 5\,(310)$).
\end{itemize}
Under the dual mapping $(m', n') = (m+n, m-n)$, $\tan\theta_{m',n'} = \tan(45^\circ - \theta_{mn})$, so $\Delta\theta_{m',n'} = 90^\circ - \Delta\theta_{mn}$ and $\Sigma(m',n') = \Sigma(m,n)$.
\end{keyidea}

On the 2D triangular lattice ($\mathbf{a}_1 = a_0(1,0)\T$, $\mathbf{a}_2 = a_0(1/2, \sqrt{3}/2)\T$), the corresponding coprime generator law for $\gcd(m,n) = 1$ ($1 \le n < m$) is:
\begin{equation}
  \tan\theta_{mn}^{\text{tri}} = \frac{\sqrt{3}\,n}{2m + n}, \qquad
  \Delta\theta_{mn}^{\text{tri}} = 2\arctan\!\left(\frac{\sqrt{3}\,n}{2m + n}\right), \qquad
  \Sigma_{\text{tri}}(m,n) = \frac{m^2 + mn + n^2}{\gcd(m^2 + mn + n^2,\, 3)}.
\end{equation}

\subsection{Structural Period $L_{\text{CSL}}$, Layer Spacing $d_{mn}$, and Commensurate Sizing}
The structural repeat period along $\Gamma_{\text{GB}}$ and the interplanar spacing of atomic layers parallel to $\Gamma_{\text{GB}}$ are:
\begin{equation}
  L_{\text{CSL}}(m,n) = a_0\sqrt{m^2 + n^2}, \qquad
  d_{mn} = \frac{a_0^2}{L_{\text{CSL}}(m,n)} = \frac{a_0}{\sqrt{m^2 + n^2}}.
\end{equation}
To ensure an integer number $K_{\text{CSL}} \ge 2$ of full structural units along the boundary length, the specimen grid size $N^{\text{CSL}}$ around a user target $N_{\text{target}}$ is quantized as
\begin{equation}
  K_{\text{CSL}} = \max\!\left(2,\; \operatorname{round}\!\left(\frac{N_{\text{target}}\,a_0}{L_{\text{CSL}}(m,n)}\right)\right), \qquad
  N^{\text{CSL}}(m,n) = \operatorname{round}\!\left(K_{\text{CSL}}\sqrt{m^2 + n^2}\right).
\end{equation}

\begin{table}[htbp]
\centering
\caption{Catalog of Special 2D Square-Lattice CSL Symmetric Tilt Grain Boundaries ($\Sigma \le 65$).}
\label{tab:csl_catalog}
\small
\begin{tabular}{ccccccccc}
\toprule
$(m,n)$ & $m^2+n^2$ & $\Sigma$ & Plane $(hkl)$ & $\theta_{mn}$ & $\Delta\theta_{mn} = 2\theta_{mn}$ & $L_{\text{CSL}}/a_0$ & $d_{mn}/a_0$ & $\delta_\perp(m,n)/a_0$ \\
\midrule
$(1,0)$ & $1$  & $\Sigma 1$   & $(100)$ & $0.000^\circ$  & $0.000^\circ$  & $1.0000$ & $1.0000$ & $0.0000$ \\
$(9,1)$ & $82$ & $\Sigma 41\text{a}$ & $(910)$ & $6.340^\circ$  & $12.680^\circ$ & $9.0554$ & $0.1104$ & $0.2774$ \\
$(7,1)$ & $50$ & $\Sigma 25\text{a}$ & $(710)$ & $8.130^\circ$  & $16.260^\circ$ & $7.0711$ & $0.1414$ & $0.3743$ \\
$(6,1)$ & $37$ & $\Sigma 37\text{a}$ & $(610)$ & $9.462^\circ$  & $18.925^\circ$ & $6.0828$ & $0.1644$ & $0.2824$ \\
$(5,1)$ & $26$ & $\Sigma 13\text{a}$ & $(510)$ & $11.310^\circ$ & $22.620^\circ$ & $5.0990$ & $0.1961$ & $0.1555$ \\
$(4,1)$ & $17$ & $\Sigma 17\text{a}$ & $(410)$ & $14.036^\circ$ & $28.072^\circ$ & $4.1231$ & $0.2425$ & $0.4549$ \\
$(7,2)$ & $53$ & $\Sigma 53\text{a}$ & $(720)$ & $15.945^\circ$ & $31.891^\circ$ & $7.2801$ & $0.1374$ & $0.3906$ \\
$(3,1)$ & $10$ & $\Sigma 5\text{a}$  & $(310)$ & $18.435^\circ$ & $36.870^\circ$ & $3.1623$ & $0.3162$ & $0.3075$ \\
$(5,2)$ & $29$ & $\Sigma 29\text{a}$ & $(520)$ & $21.801^\circ$ & $43.603^\circ$ & $5.3852$ & $0.1857$ & $0.1972$ \\
$(2,1)$ & $5$  & $\Sigma 5\text{b}$  & $(210)$ & $26.565^\circ$ & $53.130^\circ$ & $2.2361$ & $0.4472$ & $0.0456$ \\
$(7,4)$ & $65$ & $\Sigma 65\text{a}$ & $(740)$ & $29.745^\circ$ & $59.490^\circ$ & $8.0623$ & $0.1240$ & $0.1958$ \\
$(5,3)$ & $34$ & $\Sigma 17\text{b}$ & $(530)$ & $30.964^\circ$ & $61.928^\circ$ & $5.8310$ & $0.1715$ & $0.2540$ \\
$(3,2)$ & $13$ & $\Sigma 13\text{b}$ & $(320)$ & $33.690^\circ$ & $67.380^\circ$ & $3.6056$ & $0.2774$ & $0.3853$ \\
$(4,3)$ & $25$ & $\Sigma 25\text{b}$ & $(430)$ & $36.870^\circ$ & $73.740^\circ$ & $5.0000$ & $0.2000$ & $0.1400$ \\
$(5,4)$ & $41$ & $\Sigma 41\text{b}$ & $(540)$ & $38.660^\circ$ & $77.320^\circ$ & $6.4031$ & $0.1562$ & $0.3153$ \\
$(1,1)$ & $2$  & $\Sigma 1$   & $(110)$ & $45.000^\circ$ & $90.000^\circ$ & $1.4142$ & $0.7071$ & $0.0000$ \\
\bottomrule
\end{tabular}
\end{table}

% ============================================================================
\section{Microscopic Degrees of Freedom, Seam Overlap Resolution, and Pre-Opening}
% ============================================================================

\subsection{Coincidence Capture Band and Rigid Half-Crystal Translations}
Both half-crystals are anchored at $\mathbf{p}_{\text{GB}}$ using a symmetric coincidence capture band $\epsilon_{\text{band}} = 0.02\,a_0$ so floating-point roundoff ($\sim 10^{-15}\,a_0$) never drops shared coincidence sites on $s(\mathbf{X}) = 0$:
\begin{align}
  \text{Grain 0:}\quad & \mathbf{X}_{i,j}^{(0)} = \mathbf{p}_{\text{GB}} + i\,\mathbf{R}(\theta_1)\mathbf{a}_1 + j\,\mathbf{R}(\theta_1)\mathbf{a}_2, \qquad s(\mathbf{X}_{i,j}^{(0)}) \le +\epsilon_{\text{band}}, \\
  \text{Grain 1:}\quad & \mathbf{X}_{i,j}^{(1)} = \tilde{\mathbf{X}}_{i,j}^{(1)} + \delta_\parallel a_0\,\hat{\mathbf{t}}_{\text{GB}} + \delta_\perp a_0\,\hat{\mathbf{n}}_{\text{GB}}, \qquad s(\tilde{\mathbf{X}}_{i,j}^{(1)}) \ge -\epsilon_{\text{band}},
\end{align}
where $\delta_\parallel$ (\code{gb\_shift\_tangent\_a0}) and $\delta_\perp$ (\code{gb\_normal\_gap\_a0}) are the microscopic rigid-body translations parallel and normal to $\Gamma_{\text{GB}}$.

\subsection{Distance-Sorted Spatial-Hash Seam Overlap Filter}
All cross-grain pairs $(p \in G_0,\, q \in G_1)$ with Euclidean separation $r_{pq} = \|\mathbf{X}_p - \mathbf{X}_q\|_2 < d_{\min} a_0$ (default $d_{\min} = 0.45$) are extracted via a uniform spatial hash of cell size $h = d_{\min} a_0$ and sorted in ascending order of $r_{pq}$. Under the default \code{MergeMidpoint} policy, each non-removed pair $(p,q)$ is replaced by its symmetric midpoint
\begin{equation}
  \mathbf{X}_{\text{mid}} = \frac{1}{2}\left(\mathbf{X}_p + \mathbf{X}_q\right),
\end{equation}
subject to the secondary crowding guard $\|\mathbf{X}_k - \mathbf{X}_{\text{mid}}\|_2 \ge 0.85\,d_{\min} a_0$ for all surviving third neighbors $k \notin \{p,q\}$. This places shared CSL structural-unit vertices strictly on the mirror symmetry plane of the boundary.

\subsection{Derivation of the Analytical Normal Pre-Opening Gap $\delta_\perp(m,n)$}
Across a symmetric tilt boundary with layer spacing $d_{mn} = a_0/\sqrt{m^2+n^2}$, mirror-symmetric atomic planes of index $k \in \Z_{>0}$ sit at $s = \pm k\,d_{mn}$ with cross-boundary separation $r_k = 2k\,d_{mn}$. Since all layers with $r_k < d_{\min} a_0$ are fused by the overlap filter, the first unmerged cross-boundary layer pair has index and separation:
\begin{equation}
  k_{\text{first}}(m,n) = \max\!\left(1,\; \left\lceil \frac{d_{\min}\sqrt{m^2+n^2}}{2} \right\rceil\right), \qquad
  r_{\text{first}}(m,n) = 2\,k_{\text{first}}(m,n)\,d_{mn}.
\end{equation}
When $r_{\text{first}}(m,n) < 0.94\,a_0$, the initial Delaunay triangles bridging the seam are born under severe normal compression ($\det\mathbf{F}_e \ll 1$). The analytical pre-opening gap
\begin{equation}
  \delta_\perp(m,n) = \operatorname{clamp}\!\left(\frac{0.94\,a_0 - r_{\text{first}}(m,n)}{a_0},\; 0.0,\; 0.42\right)
\end{equation}
restores the initial cross-seam bond length to $r_{\text{first}} + \delta_\perp a_0 \approx 0.94\,a_0$, placing all interfacial elements inside the quadratic basin of the Zanzotto potential before L-BFGS begins.

% ============================================================================
\section{Constrained Delaunay Triangulation, MTM Energy, and Core Surgery}
% ============================================================================

Each triangle $e \in \mathcal{T}_{\text{act}}$ (filtered by geometric cutouts and maximum chord length $l_{\max} = 2.60\,a_0$) has deformation gradient $\mathbf{F}_e = [\mathbf{x}_{i_1}-\mathbf{x}_{i_0},\, \mathbf{x}_{i_2}-\mathbf{x}_{i_0}]\mathbf{D}_{0,e}^{-1}$ and Cauchy--Green metric $\mathbf{C}_e = \mathbf{F}_e\T\mathbf{F}_e$. Under Gauss--Lagrange reduction $\mathbf{C}_{\text{red},e} = \mathbf{M}_e\T \mathbf{C}_e \mathbf{M}_e \in \mathcal{D}_{\text{red}}$ with $\mathbf{M}_e \in \mathrm{GL}(2,\Z)$, the total potential energy and element Piola--Kirchhoff stress are:
\begin{align}
  E_{\text{MTM}}\big(\{\mathbf{x}_i\},\, \mathcal{T}_{\text{act}}\big) &= \sum_{e \in \mathcal{T}_{\text{act}}} \frac{A_{0,e}}{A_{\text{cell}}} \Phi_{\text{Zanzotto}}(\mathbf{C}_{\text{red},e}) + E_{\text{indenter}}\big(\{\mathbf{x}_i\}, \alpha\big), \\
  \mathbf{P}_e &= \frac{2}{A_{\text{cell}}}\,\mathbf{F}_e\,\mathbf{M}_e\,\frac{\partial \Phi_{\text{Zanzotto}}}{\partial \mathbf{C}_{\text{red}}}\,\mathbf{M}_e\T.
\end{align}
During energy-driven Delaunay surgery (\code{perform\_remeshing\_loop\_reduction}), a candidate retriangulation $\mathcal{T}_{\text{cand}} = \operatorname{Delaunay}(\mathbf{x}_{\text{eq}})$ is accepted if and only if $\Delta N_{\text{tris}} > 0$, $\mathcal{T}_{\text{cand}}$ has not been visited in a cycle, and the post-flip relaxed energy strictly decreases:
\begin{equation}
  \Delta E_{\text{surgery}} = E(\mathbf{x}_{\text{trial}}, \mathcal{T}_{\text{cand}}) - E(\mathbf{x}_{\text{old}}, \mathcal{T}_k) < -10^{-6}\,|E(\mathbf{x}_{\text{old}}, \mathcal{T}_k)|.
\end{equation}

% ============================================================================
\section{Mathematics of Associated Boundary Conditions (Relaxation vs.\ Loading)}
% ============================================================================

\subsection{5-State Independent $(x,y)$ Kinematic Classification}
Each node $i \in \{1,\dots,N\}$ accumulates boolean locks $(\ell_{x,i}, \ell_{y,i}) \in \{0,1\}^2$, drive flags $(d_{x,i}, d_{y,i}) \in \{0,1\}^2$, and velocity $\mathbf{v}_i \in \R^2$ across all matching \code{BoundaryZoneRule} entries, mapping to the 5-state role \code{node\_bc\_role[i]}:
\begin{equation}
  \text{role}(i) =
  \begin{cases}
    0 \quad (\text{Free}): & \ell_{x,i} = 0 \;\land\; \ell_{y,i} = 0 \qquad (\text{both } u_{x,i}, u_{y,i} \text{ free}), \\
    1 \quad (\text{Fixed}): & \ell_{x,i} = 1 \;\land\; \ell_{y,i} = 1 \;\land\; \neg(d_{x,i} \lor d_{y,i}) \qquad (u_{x,i} = u_{y,i} = 0), \\
    2 \quad (\text{Driven}): & \ell_{x,i} = 1 \;\land\; \ell_{y,i} = 1 \;\land\; (d_{x,i} \lor d_{y,i}) \qquad (\mathbf{u}_i(\alpha) = \alpha\mathbf{v}_i), \\
    3 \quad (\text{RollerY / X-Locked}): & \ell_{x,i} = 1 \;\land\; \ell_{y,i} = 0 \qquad (u_{x,i}(\alpha) = \alpha v_{x,i},\; u_{y,i} \text{ free}), \\
    4 \quad (\text{RollerX / Y-Locked}): & \ell_{x,i} = 0 \;\land\; \ell_{y,i} = 1 \qquad (u_{x,i} \text{ free},\; u_{y,i}(\alpha) = \alpha v_{y,i}).
  \end{cases}
\end{equation}
At any corner where a \code{RollerY} edge ($\ell_x=1$) intersects a \code{RollerX} edge ($\ell_y=1$), boolean composition yields $\ell_x=\ell_y=1$ (\code{role} $= 1$, Fixed), pinning the corner automatically.

\subsection{Subspace Projection Operator $\mathbf{P}_{\text{BC}}$ and Preconditioned Hessian Restriction}
All nodes with $\text{role}(i) \in \{0,3,4\}$ enter the active solver vector $\mathbf{z} \in \R^{2n_{\text{vars}}}$ with $z_k = x_{i(k)}$ and $z_{n_{\text{vars}}+k} = y_{i(k)}$ ($k = 0,\dots,n_{\text{vars}}-1$). The diagonal orthogonal projector $\mathbf{P}_{\text{BC}} = \operatorname{diag}(p_0, \dots, p_{2n_{\text{vars}}-1})$ zeros the $x$-component when $\text{role}=3$ and the $y$-component when $\text{role}=4$. Restricting the sparse CHOLMOD stiffness preconditioner $\mathbf{H}$ to the admissible subspace $\mathcal{V}_{\text{BC}} = \operatorname{im}(\mathbf{P}_{\text{BC}})$ via
\begin{equation}
  \tilde{\mathbf{H}} = \mathbf{P}_{\text{BC}}\,\mathbf{H}\,\mathbf{P}_{\text{BC}} + (\mathbf{I} - \mathbf{P}_{\text{BC}})
\end{equation}
guarantees that the preconditioned L-BFGS search direction $\mathbf{d} = -\tilde{\mathbf{H}}^{-1}(\mathbf{P}_{\text{BC}}\nabla E)$ satisfies $(\mathbf{I} - \mathbf{P}_{\text{BC}})\mathbf{d} = \mathbf{0}$ identically at every iteration.

\subsection{Proof of the Stage 1 Relaxation Step BCs (\code{GB-Frame Rollers})}
For a vertical bicrystal ($\phi_{\text{GB}} = 0^\circ$), the default \textbf{Relaxation Step BCs} assign \code{Left = RollerY} ($u_x=0$), \code{Right = FreeSurface}, and \code{Top/Bottom = RollerX} ($u_y=0$):
\begin{enumerate}
  \item \textbf{Rigid-Body Mode Elimination}: Any 2D rigid-body field $\mathbf{u}(x,y) = (c_x - \omega y,\, c_y + \omega x)\T$ satisfying $u_x(x_{\min}, y) = 0$ and $u_y(x, y_{\min}) = 0$ must have $c_x = c_y = \omega = 0$.
  \item \textbf{Unconstrained Normal Breathing ($\sigma_{xx}^{\text{eq}} = 0$)}: Because the Right Edge is traction-free and the Top/Bottom edges carry zero horizontal shear traction ($P_{xy}=0$ since $u_x$ is free), global horizontal equilibrium enforces $\langle P_{xx} \rangle_\Omega = 0$. The right half-crystal translates freely along $x$ to accommodate the exact equilibrium GB excess volume $\Delta u_\perp^{\text{eq}}$.
  \item \textbf{End Wedge Suppression}: Locking $u_y = 0$ on Top and Bottom enforces mirror symmetry across $y = y_{\min}, y_{\max}$ while letting $u_x$ slide freely without Poisson end-pinning.
\end{enumerate}

% ============================================================================
\section{Two-Stage Handover ($\mathbf{X}_0 \leftarrow \mathbf{x}_{\text{eq}}$) and GB Excess Energy $\gamma_{\text{GB}}(\Delta\theta)$}
% ============================================================================

At the transition from Stage 1 (\textbf{Relaxation Step BCs}) to Stage 2 (\textbf{Loading Step BCs \& Indenter}), the solver executes the reference handover operator:
\begin{equation}
  \mathbf{X}_{0,i} \longleftarrow \mathbf{x}_{\text{eq},i} \qquad \forall i \in \{1,\dots,N\},
\end{equation}
while keeping the crystallographic element reference matrices $\mathbf{D}_{0,e}$ unchanged. Thus Stage 2 trajectories $\mathbf{x}_i^{\text{Stage 2}}(\alpha) = \mathbf{x}_{\text{eq},i} + \alpha\mathbf{v}_i$ start continuously from $\mathbf{x}_{\text{eq}}$ at $\alpha = 0$ without boundary snap ($\Delta\mathbf{u}_{\text{snap}} = \mathbf{0}$), and Stage 2 nanoindentation re-anchors its contact elevation $y_{\text{top}}^{\text{eq}}$ directly on the relaxed top surface.

Finally, excluding an outer boundary collar of width $w_{\text{collar}} = 4\,a_0$ defines the interior gauge region $\Omega_{\text{int}}$ and the collar-excluded specific grain boundary excess energy:
\begin{equation}
  \gamma_{\text{GB}}^{\text{int}}(\Delta\theta) = \frac{1}{L_{\text{GB}}^{\text{int}}} \sum_{\substack{e \in \mathcal{T}_{\text{eq}} \\ \mathbf{x}_{\text{centroid},e} \in \Omega_{\text{int}}}} \frac{A_{0,e}}{A_{\text{cell}}} \Phi_{\text{Zanzotto}}(\mathbf{C}_{\text{red},e}).
\end{equation}

\end{document}
